% Original intervention schematic. Dashed arrows are deleted incoming edges.
\begin{tikzpicture}[x=1mm,y=1mm,font=\small,
  pgm variable/.style={circle,draw=BookInk,fill=BookPaper,
    minimum size=7mm,inner sep=0pt},
  target/.style={pgm variable,draw=BookAmber,fill=BookAmber!12!white},
  edge/.style={-{Stealth[length=2mm]},draw=BookInk,thick},
  removed/.style={edge,draw=BookMuted,dashed},
  heading/.style={font=\heiti\footnotesize,text=BookInk},
  result/.style={font=\footnotesize,text=BookTeal,align=center}]
  \node[heading] at (31,75) {$\operatorname{do}(X_2=1)$};
  \node[heading,align=center] at (89,75)
    {$p(X_1=1)$\\$p(X_3=1)$};

  \node[heading] at (31,63) {正向链};
  \node[pgm variable] (a1) at (7,51) {$X_1$};
  \node[target] (a2) at (31,51) {$X_2$};
  \node[pgm variable] (a3) at (55,51) {$X_3$};
  \draw[removed] (a1)--(a2);
  \draw[edge] (a2)--(a3);
  \node[result] at (89,51) {$0.5$\\$0.8$};

  \node[heading] at (31,38) {分叉};
  \node[pgm variable] (b1) at (7,26) {$X_1$};
  \node[target] (b2) at (31,26) {$X_2$};
  \node[pgm variable] (b3) at (55,26) {$X_3$};
  \draw[edge] (b2)--(b1);
  \draw[edge] (b2)--(b3);
  \node[result] at (89,26) {$0.9$\\$0.8$};

  \node[heading] at (31,13) {反向链};
  \node[pgm variable] (c1) at (7,1) {$X_1$};
  \node[target] (c2) at (31,1) {$X_2$};
  \node[pgm variable] (c3) at (55,1) {$X_3$};
  \draw[edge] (c2)--(c1);
  \draw[removed] (c3)--(c2);
  \node[result] at (89,1) {$0.9$\\$0.5$};

  \draw[BookRule] (69,-6)--(69,66);
  \node[font=\footnotesize,text=BookMuted] at (52,-14)
    {赭色节点：干预目标；虚线：删除的入边};
\end{tikzpicture}
